\clearpage
\VerdierPhysicalPageAnchor{0036}%
\VerdierSemanticAnchor{ch1-1-1-2-degree-e-morphism-cont}
其次数为 $e$，并满足通常的相容性质。

\VerdierSemanticAnchor{ch1-1-1-3}
\noindent\textbf{定义 1.1.3.}\quad
若一个 $G$-分次范畴的纤维范畴 $\mathcal C_e$ 是加性的 [20]，就称该
$G$-分次范畴是\emph{加性的}。
\VerdierSemanticAnchor{ch1-1-1-3-additive-functor}
加性分次范畴之间的一个分次函子，若其在纤维上的限制是加性函子，就称它是
\emph{加性的}。

\VerdierSemanticAnchor{ch1-1-1-3-cadd}
设 $\mathcal C$ 为一个加性 $G$-分次范畴。当 $G\neq\{e\}$ 时，忘掉纤维化结构
后的范畴 $\mathcal C$ \emph{并不是加性范畴}。不过，可以如下地给 $\mathcal C$
关联一个加性范畴：对 $\mathcal C$ 的每一对对象 $X,Y$ 和 $G$ 的每个元素 $g$，
集合 $\operatorname{Hom}_{\mathcal C}^{g}(X,Y)$ 自然带有阿贝尔群结构。定义一个
新范畴 $\mathcal C_{\mathrm{add}}$：以 $\mathcal C$ 的对象作为其对象，并把对象
$X$ 与 $Y$ 之间的态射集定义为：

\VerdierSemanticAnchor{ch1-1-1-3-eq1}
\begingroup
\makeatletter
\tagsleft@true
\let\veqno\@@leqno
\makeatother
\begin{equation}
  \operatorname{Hom}_{\mathcal C_{\mathrm{add}}}(X,Y)
  =\bigoplus_{g\in G}\operatorname{Hom}_{\mathcal C}^{g}(X,Y).
  \tag{1.1.3.1}
\end{equation}
\endgroup

\VerdierSemanticAnchor{ch1-1-1-3-cadd-composition}
$\mathcal C_{\mathrm{add}}$ 中态射的复合以显然的方式定义；可以验证
$\mathcal C_{\mathrm{add}}$ 是一个加性范畴。
\VerdierSemanticAnchor{ch1-1-1-3-fadd}
同样，给定两个加性分次范畴之间的一个加性笛卡尔函子
$F\colon\mathcal C\to\mathcal C'$，同一方法构造出一个加性函子
$F_{\mathrm{add}}\colon\mathcal C_{\mathrm{add}}\to
\mathcal C'_{\mathrm{add}}$。我们将从不使用这一构造。

\VerdierSemanticAnchor{ch1-1-1-3-restrictive}
请注意，这里采用的分次范畴定义比自然想到的定义更为严格。除了显然的结构
(1.1.2.1) 与 (1.1.2.2) 之外，我们还要求性质 (1.1.2.3)。

\VerdierSemanticAnchor{ch1-1-2}
\subsection{群在范畴上的同构意义下作用}

\VerdierSemanticAnchor{ch1-1-2-1}
\noindent\textbf{1.2.1.}\quad
设 $\mathcal C$ 为一个 $G$-分次范畴。为纤维化范畴
$\mathcal C\xrightarrow{\deg}\mathcal G$ 选取一个正规劈分 [18]，也就是说，
对每个 $g\in G$ 选取一个基变换函子：

\VerdierSemanticAnchor{ch1-1-2-1-eq1}
\begingroup
\makeatletter
\tagsleft@true
\let\veqno\@@leqno
\makeatother
\begin{equation}
  \begin{cases}
    T(g)\colon\mathcal C_e\longrightarrow\mathcal C_e,\\[0.4em]
    T(e)=\text{$\mathcal C_e$ 上的恒等函子。}
  \end{cases}
  \tag{1.2.1.1}
\end{equation}
\endgroup

\VerdierSemanticAnchor{ch1-1-2-1-transition}
对 $G$ 中每一对元素 $(g,h)$，都有 [18] 一个函子同构：

\VerdierSemanticAnchor{ch1-1-2-1-eq2}
\begingroup
\makeatletter
\tagsleft@true
\let\veqno\@@leqno
\makeatother
\begin{equation}
  \begin{cases}
    c(h,g)\colon T(hg)\longrightarrow T(g)T(h),\\[0.4em]
    c(g,e)=c(e,g)=\text{$T(g)$ 上的恒等同构}
  \end{cases}
  \tag{1.2.1.2}
\end{equation}
\endgroup
